% !PRINTED 150
\weberpage{150}
Die Differenz
\[
S - P = S'
\]
ist also wieder eine symmetrische Function, deren h\"ochstes Glied
niedriger ist als das von $S$, und die wir nach der Voraus\-setzung
rational durch die Grundfunctionen darstellen k\"onnen.
Da $P$ ebenso dargestellt ist, so ist das Ziel erreicht. Es erhellt
auch hier wieder, dass durch die Darstellung mittelst der Grund\-functionen
keine Br\"uche eingef\"uhrt werden, wenn urspr\"unglich
in $S$ keine enthalten sind.

Diese Methode der Berechnung symmetrischer Functionen
giebt uns zugleich Aufschluss \"uber den Grad der ganzen ratio\-nalen
Function von $a_1,\ a_2 \ldots a_n$, durch die eine bestimmte
symmetrische Function der $\alpha_1,\ \alpha_2 \ldots \alpha_n$ ausgedr\"uckt wird.
Wenn n\"amlich $(\nu_1,\ \nu_2 \ldots \nu_n)$ die Ordnung der symmetrischen
Function ist (bemessen nach der Exponentenreihe des ersten
Gliedes), so ist $\nu_1$ der Grad von $P$ in Bezug auf s\"ammtliche
$a_1, a_2 \ldots a_n$, und in dem entwickelten Ausdruck f\"ur $S$ kann kein
Glied h\"oheren Grades vorkommen; auch kann das Glied $P$ durch
keines der folgenden Glieder zerst\"ort werden, da $P$ durch die
Ordnung von $S$ v\"ollig bestimmt ist, und die Ordnung von $S'$
und aller folgenden symmetrischen Functionen niedriger ist, als
die Ordnung von $S$.

Setzen wir also an Stelle von $a_1, a_2 \ldots a_n$:
\[
\frac{a_1}{a_0},\quad \frac{a_2}{a_0} \ldots \frac{a_n}{a_0},
\]
% [corrected 2026-08-06, 12x zoom strip 8/12] The letterspacing starts at "ganze", not at
% "eine": on the scan "eine" is set with normal spacing and every following word to "von" is
% letterspaced. Raised as a low-confidence query by this page's translation pass and confirmed.
so ist $a_0^{\nu_1}\,S$ eine \gesperrt{ganze homogene Function}
$\nu_1^{\mathrm{ten}}$ \gesperrt{Grades von} $a_0,\ a_1 \ldots a_n$, und man kann $S$ nicht durch Multiplication
mit einer niedrigeren Potenz von $a_0$ in eine ganze Function der
$a$ verwandeln.

\begin{center}
\S.\,46.

\gesperrt{Discriminanten.}
\end{center}

Die im Vorhergehenden gewonnenen S\"atze lehren, sym\-metrische
Functionen der Wurzeln einer Gleichung rational durch
die Co\"efficienten der Gleichung auszudr\"ucken; die Resultate sind
zun\"achst abgeleitet unter der Voraussetzung, dass die Wurzeln,
oder, was nach dem dritten Abschnitt damit gleichbedeutend ist,
%CONT>>
